% COL774 Assignment 1.2 -- report  (Part 2: Logistic Regression)
%
% Compile: pdflatex report.tex   (geometry, booktabs, amsmath, graphicx)
% The four PNGs must sit next to this file.
% Submit as: 2023CS10058_2023CS10552.zip
%     python dev/make_submission.py 2023CS10058 2023CS10552
%
% All numbers are real measurements: Parts (a)/(b) on part_ab_train/val.csv,
% Part (c) on the Kaggle competition data with the baseline model (part_c.py).
\documentclass[10pt,a4paper]{article}
\usepackage[margin=1.5cm]{geometry}
\usepackage{booktabs}
\usepackage{amsmath}
\usepackage[T1]{fontenc}
\usepackage{graphicx}

\setlength{\parskip}{2pt}
\setlength{\textfloatsep}{5pt}
\setlength{\intextsep}{5pt}
\setlength{\parindent}{0pt}
\newcommand{\feat}[1]{\texttt{\small #1}}
\pagestyle{empty}

\begin{document}

\begin{center}
  {\large\bfseries COL774 Assignment 1.2 --- Report}\\[1pt]
  {\small Popat Nihal Alkesh (2023CS10058) \qquad Viraaj Narolia (2023CS10552)}
\end{center}
\vspace{-2pt}\hrule\vspace{6pt}

%======================================================================
\section*{Part (a) --- Gradient-descent variants}

Loss is plotted against \emph{wall-clock training time}, not epochs: full-batch
takes one update per epoch and SGD takes $n$, so equal epochs are not equal work.

\begin{center}
\includegraphics[width=0.47\textwidth]{parta_train_loss_vs_time.png}
\hfill
\includegraphics[width=0.47\textwidth]{parta_val_loss_vs_time.png}
\end{center}

\begin{table}[h]
\centering\footnotesize
\begin{tabular}{@{}lrrrrrr@{}}
\toprule
\textbf{Method} & \textbf{Epochs} & \textbf{Train s} & \textbf{Final train}
& \textbf{Final val} & \textbf{Min val} & \textbf{at epoch} \\
\midrule
Full-batch GD      & 500 & 0.59 & 0.4793 & 0.5339 & 0.5339 & 500 \\
Mini-batch GD      & 200 & 0.64 & 0.4652 & 0.5301 & 0.5281 & 72  \\
SGD                & 30  & 2.22 & 0.4857 & 0.5350 & 0.5350 & 30  \\
Mini-batch AdaGrad & 200 & 0.76 & \textbf{0.4626} & \textbf{0.5286} & \textbf{0.5265} & 193 \\
\bottomrule
\end{tabular}
\end{table}

\textbf{Fastest.} Mini-batch and AdaGrad reach $0.60$ validation loss in about
$10$\,ms; full-batch needs $\sim0.2$\,s and SGD $\sim0.5$\,s. One pass buys
full-batch a single step but buys mini-batch $n/32 \approx 135$, so it does the
same arithmetic to move the parameters far less. AdaGrad matches mini-batch early
and ends lowest on both losses, its per-coordinate scaling paying off late once
the raw gradient has shrunk. SGD is slowest in wall-clock terms despite the
fewest epochs: at batch size $1$ it runs $n$ Python-level iterations per epoch,
so per-update overhead dominates. The three batched methods' totals differ by
less than timing noise, so the curve shape is the meaningful comparison.

\textbf{Overfitting.} Mini-batch bottoms out at epoch $72$ of $200$ and rises
after, so it overfits for two thirds of its schedule. AdaGrad bottoms at epoch
$193$ of $200$ --- still improving when the schedule ends. Full-batch and SGD
never turn: their minima are their final epochs, so within the prescribed budgets
they are under-trained rather than over-trained.

%======================================================================
\section*{Part (b) --- Class imbalance}

Counts are N${=}2723$, AF${=}366$, Other${=}1221$ ($\approx 7{:}1$ against AF).
All four models share the fixed AdaGrad protocol, so differences are due to the
reweighting alone. Metrics on \texttt{part\_ab\_val.csv}:

\begin{table}[h]
\centering\footnotesize
\begin{tabular}{@{}lrrrrrrr@{}}
\toprule
\textbf{Method} & \textbf{Macro-AP} & \textbf{Bal.\ acc} & \textbf{Plain acc}
& \textbf{Rec.\ N} & \textbf{Rec.\ AF} & \textbf{Rec.\ O} & \textbf{Prec.\ AF} \\
\midrule
Baseline       & 0.7765 & 0.7186 & \textbf{0.7998} & 0.9366 & 0.6923 & 0.5267 & 0.8060 \\
classweight    & 0.7748 & \textbf{0.7555} & 0.7738 & 0.8305 & \textbf{0.7949} & 0.6412 & 0.6200 \\
classweight2   & \textbf{0.7777} & 0.7322 & 0.7965 & 0.9161 & 0.7308 & 0.5496 & 0.7600 \\
Focal ($\gamma{=}2$) & 0.7659 & 0.7403 & 0.7846 & 0.8767 & 0.7564 & 0.5878 & 0.6782 \\
\bottomrule
\end{tabular}
\end{table}

\textbf{Which metrics reveal the benefit.} Two of these columns are actively
misleading. \emph{Plain accuracy gets worse} under every treatment
($0.7998 \to 0.7738$), because it is dominated by the majority class and
reweighting deliberately trades Normal recall ($0.9366 \to 0.8305$) for minority
recall. \emph{Macro-AP is flat} (within $\pm0.011$), because AP scores the
\emph{ranking} while reweighting mainly rescales logits and shifts the decision
boundary --- it changes which side of the boundary examples fall on, not their
order, so a threshold-free ranking metric is nearly blind to it.

What does show the benefit is \textbf{balanced accuracy} ($0.7186 \to 0.7555$)
and \textbf{AF recall} ($0.6923 \to 0.7949$, $+10.3$ points), both of which weight
the minority class equally. The ordering is consistent --- full-strength $>$ focal
$>$ \texttt{classweight2}, whose $\alpha^{0.3}$ compresses the weights towards $1$
and so corrects least. The cost is AF precision, $0.8060 \to 0.6200$.

%======================================================================
\section*{Part (c) --- Feature engineering for AF detection}

The metric $M = \mathrm{TPR} - 100\,\mathrm{FPR}$ makes one false positive cost as
much as $100/N$ of recall, and every choice below follows from that.

\textbf{Features and preprocessing.} $\phi(x)$ is the \textbf{392 supplied
columns} ($d = 392 < 500$), median-imputed, winsorised to the training
$[0.1\%, 99.9\%]$ quantiles, then standardised --- all statistics fit on
\texttt{train.csv} only. Winsorising matters here specifically: one wild outlier
at test time would otherwise produce a huge logit and a confident false positive.
A second pipeline adding $\sim$70 raw-signal features (Pan--Tompkins QRS
detection feeding COSEn, symbolic-dynamics entropy and Lorenz-plot $\delta$RR
occupancy; P-wave and $4$--$9$\,Hz fibrillatory-wave measures via QRST
cancellation; signal-quality descriptors) was built and tested but did not beat
the supplied columns here, which is itself the finding. Following the course
announcement, patients $45$ and $51$ were dropped from training and threshold
selection ($1413$ and $1371$ rows), never from the scored set.

\textbf{(i) Top 10 features.} Columns are standardised, so coefficients are
comparable and features are ranked by $|w_j|$.

\begin{table}[h]
\centering\footnotesize
\begin{tabular}{@{}rlr@{\hskip 16pt}rlr@{}}
\toprule
\# & \textbf{Feature} & $w_j$ & \# & \textbf{Feature} & $w_j$ \\
\midrule
1 & \feat{p\_wave\_corr\_coeff\_median}      & $-0.686$ & 6  & \feat{full\_waveform\_median}    & $-0.367$ \\
2 & \feat{rr\_sampen}                        & $+0.448$ & 7  & \feat{s\_wave\_amp}              & $-0.328$ \\
3 & \feat{pri\_approximate\_entropy}         & $-0.435$ & 8  & \feat{p\_wave\_corr\_coeff\_std} & $+0.323$ \\
4 & \feat{rri\_multiscale\_entropy}          & $+0.423$ & 9  & \feat{pnn400}                    & $-0.303$ \\
5 & \feat{p\_wave\_higuchi\_fractal\_median} & $-0.391$ & 10 & \feat{t\_wave\_eng}              & $+0.294$ \\
\bottomrule
\end{tabular}
\end{table}

This recovers the two textbook criteria for AF unprompted: four P-wave features
lead, headed by beat-to-beat P-wave reproducibility with a large negative weight
(less consistent $\Rightarrow$ more AF-like, i.e.\ the absent P wave), and three
RR-irregularity measures follow, all signed so that more irregular means more AF.

\textbf{(ii) Class imbalance, and (iv) final metrics.} Class weighting was
searched ($\{$none, \texttt{balanced}$\}$ jointly with $C$, by patient-grouped CV
scored on $M$) and \emph{rejected}: `balanced' buys recall at the cost of
specificity, the right trade for Part (b)'s balanced accuracy and the wrong one
at $100\times$ per false positive. The imbalance is instead handled by
cost-sensitive thresholding --- moving the cut from $0.5$ to $0.8887$.

\begin{table}[h]
\centering\footnotesize
\begin{tabular}{@{}lllrrrrrr@{}}
\toprule
\textbf{Model} & \textbf{Population} & \textbf{Thr.} & \textbf{M} & \textbf{TP}
& \textbf{FP} & \textbf{AF recall} & \textbf{AF FPR} & \textbf{AP} \\
\midrule
Submitted (392 cols)      & val, excl.\ pt.\ 51 & $0.8887$ & $\mathbf{0.6378}$ & $2115$ & $13$ & $0.7387$ & $0.00101$ & $0.9878$ \\
Submitted (392 cols)      & val, excl.\ pt.\ 51 & $0.5$    & $0.5931$ & $2432$ & $33$ & $0.8495$ & $0.00256$ & $0.9878$ \\
Part (a) baseline (78 cols) & val, excl.\ pt.\ 51 & $0.9504$ & $0.2339$ & $714$ & $2$ & $0.2494$ & $0.00016$ & $0.9481$ \\
Part (a) baseline (78 cols) & val, excl.\ pt.\ 51 & $0.5$    & $0.1529$ & $2150$ & $77$ & $0.7510$ & $0.00598$ & $0.9481$ \\
Submitted (392 cols)      & val, all rows       & $0.8887$ & $-2.5683$ & $2893$ & $450$ & $0.7846$ & $0.03353$ & $0.9109$ \\
Submitted (392 cols)      & val, all rows       & $0.5$    & $-3.1255$ & $3256$ & $538$ & $0.8831$ & $0.04009$ & $0.9109$ \\
\bottomrule
\end{tabular}
\end{table}

Thresholding cuts false positives $33 \to 13$ for $317$ true positives, which
this metric rates as worth it; AP is identical across thresholds because it is
threshold-free, the same insensitivity seen in Part (b).

\textbf{Against the Part (a) baseline.} Restricting $\phi(x)$ to the $78$
columns used in Parts (a)/(b) costs $0.40$ of $M$ ($0.6378 \to 0.2339$) and
$0.04$ of AP. The $78$-column model is not merely noisier --- it is forced to be
far more timid, choosing a threshold of $0.9504$ and catching only $24.9\%$ of AF
to hold false positives down, where the full representation catches $73.9\%$ at a
comparable false-positive count. The top-10 table explains why: nine of those ten
features come from the $314$ Goodfellow columns, and only \feat{rr\_sampen}
belongs to the Part (a) set. The P-wave descriptors that dominate the model
simply are not available in the Part (a) representation.

\textbf{One noisy patient.} The two populations differ by $3.2$ in $M$, and
almost all of it is patient $51$: it supplies $547$ non-AF windows of which
$437$ are called AF --- an $80\%$ false-positive rate against $0.10\%$ for the
other eleven validation patients combined, i.e.\ $437$ of all $450$ false
positives. This is not a threshold artefact: AP falls $0.9878 \to 0.9109$ when it
is included, and AP does not depend on the threshold, so the model genuinely
\emph{ranks} that patient's non-AF windows as AF-like. We report both rows rather
than only the favourable one, since $M = 0.6378$ is the score conditional on the
evaluation set containing no such patient.

\textbf{(iii) Threshold selection.} Chosen on pooled patient-grouped
out-of-fold predictions on \texttt{train.csv} plus predictions on the untouched
\texttt{val.csv}, sweeping every distinct predicted probability ($17{,}108$
candidate cuts) rather than a grid. We take not the arg-max of $M$ but the
\emph{largest} threshold within $0.01$ of it: the arg-max sits exactly where a
few negatives happen to fall below the cut, and a stricter threshold costs almost
no recall while protecting against one more negative landing above it. This gave
$\mathbf{0.8887}$, against the calibrated-optimal anchor
$\alpha P/(N + \alpha P) = 0.9633$ (from admitting a window with
$P(\mathrm{AF}) = p$ changing $M$ by $p/P - \alpha(1-p)/N$); the gap says the
fitted probabilities are mildly under-confident, which the empirical sweep does
not assume away.

\begin{center}
\includegraphics[width=0.40\textwidth]{partc_m_vs_threshold.png}
\hfill
\includegraphics[width=0.56\textwidth]{partc_epsilon_sweep.png}
\end{center}

\textbf{Trend in $1/\epsilon$} (sampled at $200$ points over $[10,1000]$). As the
penalty grows the optimal threshold rises towards $1$ and the best achievable $M$
falls: the anchor $\alpha P/(N + \alpha P)$ increases monotonically in $\alpha$,
so more confidence is demanded before calling AF, and once only the most certain
windows are admitted no further increase can remove the residual false positives.
The sweep is over the full validation set, so $M$ is already negative over most of
the range --- with $\mathrm{FPR} = 0.0335$ the penalty $\alpha\,\mathrm{FPR}$
exceeds the maximum attainable recall once $1/\epsilon \gtrsim 30$. That is the
patient-$51$ effect again, seen against the penalty rather than the threshold.

%======================================================================
\section*{Disclosure}

\textbf{AI use.} Completed with assistance from Claude (Anthropic), used for code
design, implementation and drafting this report; all outputs were reviewed,
executed and verified by us. \textbf{Open source.} No third-party code is copied
into the submitted files; only NumPy, SciPy, pandas and scikit-learn are
imported. The raw-signal features (built but not submitted) re-implement
published methods cited at their definitions: Pan--Tompkins QRS detection; Lake
\& Moorman's COSEn; Zhou et al.'s symbolic-dynamics entropy; Sarkar et al.'s
Lorenz-plot $\delta$RR analysis; zero-padding QRST cancellation. The $314$
supplied Goodfellow columns come from the \texttt{Seb-Good/ecg-features} library,
as stated in the assignment.

\end{document}
